\section{What the October 2026 evidence supports}\label{sec:C34:01}
The evidence supports a system-level account of agent competence. Model architecture and training determine a conditional policy, but context construction, tools, memory, runtime, and evaluation determine which policy is actually executed and which outcomes are credited. This is not an argument that model quality is secondary. It is an identification statement: when several layers change together, an outcome difference cannot be attributed to the model or to a named memory mechanism without additional controls.

The current frontier also changes which implementation assumptions are reasonable. FlashAttention-4, current hybrid cache managers, and low-precision architectures make hardware and model layout part of any performance claim. V-JEPA 2.1 and contemporary world-model work sharpen representation questions but do not collapse predictive quality into control validity. Context Language Models, MemoryWalker, RecMem, InMind, HyperSkill, and 2026 compaction studies expose different failure points in retained state: writing, consolidation, access, implicit relevance, influence, conditioning, and lifecycle maintenance. A survey that treats them all as larger memory misses the distinctions needed to design experiments~\cite{rewrite_fa4,rewrite_vjepa21,shao2026context,j2026memorywalker,dai2026recmem,inmind2026,xu2026hyperskill}.

The strongest practical conclusion is conditional. Reliable execution requires that decision-relevant evidence, procedural applicability, authorization, and effect status remain interpretable across the system's state transitions. Longer windows and larger stores can help, but neither establishes those properties. A persistent archive can coexist with failed retrieval. Correct retrieval can coexist with inappropriate influence. A valid procedure can coexist with revoked authority. An authorized call can coexist with an unknown external outcome. Each transition needs its own evidence.

No single benchmark or finite improvement trajectory establishes unrestricted AGI. Generality remains a profile over declared task families, resources, adaptation, and risk. Recursive self-improvement remains a set of bounded mechanisms whose executor gains, improver gains, and evaluator dependence must be separated. Negative evidence, including task-order sensitivity, evaluator false acceptance, and limited transfer, is part of the scientific account rather than an inconvenience to remove from an aggregate score.

Fixed-policy interface changes and learning across interfaces also answer different questions. Harness decision geometry separates an interface's regularized optimum from the deployed policy's mismatch under finite-state assumptions; the coding experiments use a teacher proxy rather than access to the true optimum. Harness Scaling instead tests training across structured harness combinations and held-out compositions~\cite{ch01_harnessgeometry2026,ch01_harnessscaling2026}. Together they motivate measuring compatibility and adaptation separately. Neither a mismatch decomposition nor successful transfer among sampled interfaces establishes robustness to arbitrary future tool semantics, memory transformations, or authority changes.

